\chapter{Decimale getallen}\label{ch:g5:decimals}

Tienden en honderdsten verschenen in \cref{ch:g4:decimals}; \'{e}\'{e}n
niveau erbij, de duizendsten, maakt het decimale stelsel compleet.
Dit hoofdstuk consolideert lezen, ontbinden en --- vooral ---
decimale getallen vergelijken zonder in hun beroemde valkuilen te
trappen.

\section{Tot de duizendsten}

\begin{definition}[Decimale plaatsen]\label{def:g5:decimals:places}
Na de \omterm{def:g4:decimals:point}{decimale punt} komen de tienden, de honderdsten, de
duizendsten --- elk tien keer minder waard dan de vorige:
\[
4.362 = 4 + \frac{3}{10} + \frac{6}{100} + \frac{2}{1000}
= \frac{4\,362}{1000}.
\]
\end{definition}

\begin{omfigure}
\begin{tikzpicture}[scale=0.9]
  \draw (0,0) grid[xstep=2.3, ystep=0.85] (9.2,1.7);
  \node[font=\small] at (1.15,1.275) {eenheden};
  \node[font=\small] at (3.45,1.275) {tienden};
  \node[font=\small] at (5.75,1.275) {honderdsten};
  \node[font=\small] at (8.05,1.275) {duizendsten};
  \node[font=\large, omDef] at (1.15,0.425) {$4$};
  \node[font=\large, omThm] at (3.45,0.425) {$3$};
  \node[font=\large, omThm] at (5.75,0.425) {$6$};
  \node[font=\large, omThm] at (8.05,0.425) {$2$};
  \fill (2.3,0.06) circle (2.5pt);
\end{tikzpicture}

{\small De plaatswaardetabel van $4.362$: hele deel in blauw,
decimaal deel in rood, de punt tussen eenheden en tienden.}
\end{omfigure}

\begin{example}[Veel lezingen, \'{e}\'{e}n getal]\label{ex:g5:decimals:readings}
$4.362$ kan gelezen worden als ``vier komma drie zes twee'', of
``vier en $362$ duizendsten'', of ontbonden als
$4 + 0.3 + 0.06 + 0.002$. En opvullen met eindnullen verandert
niets: $4.362 = 4.3620$. Maar $4.362 \neq 4.0362$: een \omterm{def:g1:counting:zero}{nul}
\emph{binnenin} duwt elk cijfer naar een kleinere plaats.
\end{example}

\begin{example}[Waar staat het op de lijn?]\label{ex:g5:decimals:line}
Om $4.362$ te plaatsen: tussen $4$ en $5$; inzoomen, tussen $4.3$
en $4.4$; opnieuw inzoomen, tussen $4.36$ en $4.37$, dichter bij
$4.36$. Elk decimaal cijfer is \'{e}\'{e}n zoomniveau meer.
\end{example}

\begin{omfigure}
\begin{tikzpicture}[scale=1.15]
  \draw[-Stealth] (-0.2,1.6) -- (10.4,1.6);
  \foreach \i in {0,...,10} \draw (\i,1.52) -- (\i,1.68);
  \node[below, font=\small] at (0,1.5) {$4$};
  \node[below, font=\small] at (10,1.5) {$5$};
  \node[below, font=\small] at (3,1.5) {$4.3$};
  \node[below, font=\small] at (4,1.5) {$4.4$};
  \fill[omThm] (3.62,1.6) circle (1.7pt);
  \draw[omExo, thin] (3,1.45) -- (0,0.15);
  \draw[omExo, thin] (4,1.45) -- (10,0.15);
  \draw[-Stealth] (-0.2,0) -- (10.4,0);
  \foreach \i in {0,...,10} \draw (\i,-0.08) -- (\i,0.08);
  \node[below, font=\small] at (0,-0.1) {$4.3$};
  \node[below, font=\small] at (10,-0.1) {$4.4$};
  \node[below, font=\small] at (6,-0.1) {$4.36$};
  \fill[omThm] (6.2,0) circle (1.7pt)
    node[above, font=\small] {$4.362$};
\end{tikzpicture}

{\small Inzoomen op de getallenlijn: het stuk van $4.3$ tot $4.4$,
vergroot, is zelf in tien honderdsten gesneden --- en $4.362$ zit
net voorbij $4.36$.}
\end{omfigure}

\section{Decimalen vergelijken}

\begin{method}[Twee decimalen vergelijken]\label{met:g5:decimals:compare}
\begin{enumerate}
  \item Vergelijk de hele delen: $6.1 > 5.987$.
  \item Gelijke hele delen: vergelijk de decimale delen cijfer voor
        cijfer vanaf de punt --- tienden, dan honderdsten, dan
        duizendsten;
  \item opvullen met eindnullen maakt de vergelijking eerlijk:
        $2.5$ vs $2.48$ wordt $2.50$ vs $2.48$, en $50 > 48$
        honderdsten.
\end{enumerate}
De valkuil, nog \'{e}\'{e}n keer: \emph{een langer decimaal deel betekent
niet een groter getal} --- $2.48$ heeft meer cijfers dan $2.5$ en is
kleiner.
\end{method}

\begin{example}\label{ex:g5:decimals:order}
Zet op volgorde $7.3$; $7.09$; $7.31$; $7.299$. Vul aan tot drie
decimalen: $7.300$; $7.090$; $7.310$; $7.299$. Dan
\[
7.09 < 7.299 < 7.3 < 7.31 .
\]
Merk op $7.299 < 7.3$ ook al lijkt $299$ groot: $299$ duizendsten
tegen $300$ duizendsten.
\end{example}

\begin{example}[Tussen twee decimalen wringen]\label{ex:g5:decimals:between}
Is er een getal tussen $5.7$ en $5.8$? Ja, volop: $5.75$, $5.71$,
$5.799$\,\dots\ Tussen $5.79$ en $5.8$? Weer volop: $5.795$,
bijvoorbeeld. Tussen twee verschillende decimalen kun je
\emph{altijd} nog \'{e}\'{e}n wringen --- er is geen ``volgend''
decimaal getal.
\end{example}

\section{Decimalen afronden}

\begin{definition}[Afronden]\label{def:g5:decimals:round}
Om af te ronden op de eenheid (of tiende, of honderdste), kijk naar
het volgende cijfer: $5$ of meer rondt naar boven af, anders naar
beneden. Dus $4.362$ rondt af op $4$ (eenheid), op $4.4$ (tiende),
op $4.36$ (honderdste).
\end{definition}

\begin{example}[Geld rondt af op centen]\label{ex:g5:decimals:money}
Een berekende prijs van $7.4983$ wordt getoond als $7.50$: bedragen
in het echte leven worden afgerond op de honderdste. Let op de
cascade: $7.4983 \to 7.50$, waar de $49$ $50$ werd --- \omterm{def:g4:numbers:round}{afronden} kan
door meerdere cijfers rimpelen.
\end{example}

\section{Oefeningen}

\begin{exercise}[$\star$]\label{exo:g5:decimals:1}
Schrijf als decimaal getal: $5 + \frac{2}{10} + \frac{7}{1000}$;
\quad $\frac{85}{100}$; \quad $\frac{4\,507}{1000}$; \quad ``twaalf
en negen duizendsten''.
\end{exercise}

\begin{exercise}[$\star$]\label{exo:g5:decimals:2}
In $27.418$: wat telt de $4$? De $8$? Ontbind het getal zoals in
\cref{def:g5:decimals:places}.
\end{exercise}

\begin{exercise}[$\star$]\label{exo:g5:decimals:3}
Waar of niet waar? Leg uit met plaatsen.
$3.60 = 3.6$; \quad $0.5 = 0.05$; \quad $2.070 = 2.07$; \quad
$1.3 = 1.03$.
\end{exercise}

\begin{exercise}[$\star$]\label{exo:g5:decimals:4}
Schrijf over en vul in met $<$, $>$ of $=$:
\[
4.8 \;?\; 4.79, \qquad
0.302 \;?\; 0.32, \qquad
6.5 \;?\; 6.500, \qquad
9.09 \;?\; 9.1 .
\]
\end{exercise}

\begin{exercise}[$\star$]\label{exo:g5:decimals:5}
Zet op volgorde van klein naar groot: $3.14$; \ $3.4$; \ $3.041$; \
$3.104$; \ $3.41$.
\end{exercise}

\begin{exercise}[$\star$]\label{exo:g5:decimals:6}
Tussen welke twee hele getallen ligt elk getal? Tussen welke twee
getallen met \'{e}\'{e}n decimaal? \quad $6.83$; \quad $0.492$; \quad
$12.06$.
\end{exercise}

\begin{exercise}[$\star$]\label{exo:g5:decimals:7}
Rond $8.276$ af: op de eenheid; op de tiende; op de honderdste.
Rond $3.95$ af op de tiende.
\end{exercise}

\begin{exercise}[$\star$]\label{exo:g5:decimals:8}
Geef een getal strikt tussen: $2.6$ en $2.7$; \quad $0.41$ en
$0.42$; \quad $5.99$ en $6$.
\end{exercise}

\begin{exercise}[$\star$]\label{exo:g5:decimals:9}
De hoogtes van vier planten zijn $0.85$ m, $1.2$ m, $0.9$ m en
$1.02$ m. Zet ze op volgorde van kort naar lang.
\end{exercise}

\begin{exercise}[$\star\star$]\label{exo:g5:decimals:10}
Lucie zegt: ``$7.12 > 7.9$ omdat $12 > 9$''. Corrigeer haar fout,
vergelijk eerst de tienden.
\end{exercise}

\begin{exercise}[$\star\star$]\label{exo:g5:decimals:11}
Zoek alle getallen met precies twee decimale cijfers die \omterm{def:g4:numbers:round}{afronden} op
$6.4$ bij \omterm{def:g4:numbers:round}{afronden} op de tiende. Wat zijn het kleinste en het
grootste?
\end{exercise}
